% Appendix I -- PREPRINT ONLY (\ifarxiv). Two jobs, both of which the venue version
% has no page budget for:
%
%  1. Consolidate the evidence that the model is sound WITHOUT our owning a
%     Grace-Hopper part, by checking it against other people's measurements on one.
%  2. State exactly what we would measure given access, in enough detail that
%     someone with the hardware could run it or lend it. This is the paper's ask.
%
% Written to be readable on its own by someone deciding whether to grant access.

\section{Validating the model without the hardware}
\label{app:validation}

Every positive number in this paper is a projection, and the machine that would
test it is one we do not have. That is a real limitation and we do not soften it.
What we can do is check the projection's \emph{inputs} against measurements other
people have published on the hardware in question, and be explicit about which
inputs remain unchecked. This appendix does that, and then says what one week of
access would settle.

\subsection{What the model needs, and where each input comes from}

\begin{center}\small
\begin{tabular}{llll}
\toprule
Input & Value & Source & Status \\
\midrule
C2C bandwidth $\Xlink$ & 419 GB/s & Pie, measured~\cite{pie} & external, optimistic \\
                       & 375 GB/s & Schieffer et al.~\cite{schieffer2024} & external, central \\
HBM bandwidth & 4 TB/s & vendor & spec \\
HBM efficiency $\eta_w$ & 0.75 & our A100 fits & \textbf{checked, \S\ref{app:validation:t0}} \\
$\mathrm{ITL}_0$ & 27 ms (70B) & BOOST, measured~\cite{boost} & \textbf{checked to 1.1\%} \\
the knee $r^{*}$ & Eq.~\ref{eq:knee} & BOOST derive it too~\cite{boost} & \textbf{independently derived} \\
$F$, fixed cost & $0.053$ ms/layer & our A100 ablations & \textbf{unchecked on C2C} \\
$c$, per rider & $0.177 + 0.0054L$ & our A100 ablations & \textbf{unchecked on C2C} \\
\bottomrule
\end{tabular}
\end{center}

Four of the six quantities that decide the verdict are now anchored to someone
else's Grace-Hopper measurement. The two that are not are the two we introduced.

\subsection{The step time, checked against a published GH200 measurement}
\label{app:validation:t0}

BOOST~\cite{boost} reports a GH200 decode step for a 70B-class model, which is the
measurement we could not take: Llama-3.3 70B at $1$K/$8$K context, HBM-only, batch
$6$ in $81.6$\,GB of HBM, \textbf{TPOT $= 27.1$\,ms}. Reconstructing that
configuration --- FP8 weights, FP16 KV, mean sequence $5$K --- recovers their stated
capacity to $1.5\%$, which makes the KV-precision inference self-consistent, and our
expression for $\mathrm{ITL}_0$ returns $\mathbf{26.8}$\,\textbf{ms}: an error of
$1.1\%$, with an implied HBM efficiency of $0.742$ against the $0.75$ we had assumed
from A100 fits. Their table does not state the KV dtype, so this is a cross-check
rather than a validation, and we report it as one.

It also disciplines a choice we had made for convenience. The weights-only form of
$\mathrm{ITL}_0$ returns $23.5$\,ms against the same measurement, $13\%$ low, so
omitting the resident cohort's KV read understates the window a transfer can hide
inside, and every rider count computed from it is conservative by about that much.

\subsection{The knee, derived independently by others}

The crossover of Eq.~\ref{eq:knee} is not only ours. BOOST derives the optimal share
of accesses to serve from host memory as $\alpha/(1{+}\alpha)$ for a host-to-HBM
bandwidth ratio $\alpha$, and \emph{measures} the optimum on a GH200 at $8.3\%$
against $9.1\%$ predicted. DAK~\cite{dak} reaches the same expression as
$B_h/(B_h + B_g)$. Those are the per-byte form of our per-request $r^{*}$: multiply
an optimal host \emph{fraction} by the step's total bytes and divide by the bytes one
request contributes and the two coincide. Two independent derivations and one direct
measurement of the same crossover is stronger support than we could have produced
alone, and the $8.3$-versus-$9.1$ gap is a useful warning that the closed form runs
about $10\%$ optimistic in practice.

\subsection{Which direction the unchecked inputs err}

Three of the remaining uncertainties move the verdict in our favour, which is worth
stating plainly because the reader should not have to take the projection on trust.

\begin{itemize}\itemsep1pt \parskip0pt
\item \textbf{HBM efficiency.} If the real $\eta_w$ is below $0.75$,
  $\mathrm{ITL}_0$ rises, which both raises $r^{*}$ and shrinks
  $F/\mathrm{ITL}_0$. At DirectKV's measured $1.9$ of $4$\,TB/s the projected gain
  rises from $\times 1.13$ to $\times 1.43$.
\item \textbf{The omitted KV term.} As above, $13\%$ conservative.
\item \textbf{$F$ on a Grace host.} $F$ is launch- and host-bound, and Grace's cores
  are not the A100 host's. We cannot sign this one, and say so below.
\end{itemize}

And one errs against us: \textbf{$\Xlink$}. We use Pie's $419$\,GB/s, which is the
top of the published range; Schieffer et al.\ report $375$~\cite{schieffer2024} and
BOOST $350$\,GiB/s. Since $r^{*}$ scales with $\Xlink$, a central estimate of
$375$\,GB/s would lower every rider count in this paper by about $10\%$. We keep
$419$ for continuity with the measurement it came from and flag the sensitivity
here rather than burying it.

\subsection{What one week of access would settle}
\label{app:validation:ask}

The experiments below need \emph{one} Grace-Hopper node and no new code. They are
ordered by how much they would change the paper.

\begin{enumerate}\itemsep2pt
\item \textbf{Is $F$ launch-shaped or byte-shaped?} This is the crux and it is
  cheap. $F$ is the one coefficient with no external anchor, and whether it carries
  to a coherent link depends entirely on its shape: DirectKV's PCIe-to-C2C
  comparison shows a \emph{byte}-shaped overhead halving with bandwidth
  ($10\% \to 5\%$), while a launch-shaped one should not move at all. Sweeping
  bytes per step at a fixed launch count separates them, and we can run half of
  this on an A100 today. One hour.
\item \textbf{$F$ and $c$ on Grace.} Re-run the ablation ladder of
  Appendix~\ref{app:ridercost} --- three baseline arms bracketing each ablation,
  purity-gated --- and compare the coefficients to the A100 values in
  Eq.~\ref{eq:coeffs}. Half a day, and it is the measurement the whole projection
  rests on.
\item \textbf{The knee, and the sawtooth.} Sweep context across an integer crossing
  of Eq.~\ref{eq:knee} and confirm the admitted rider count steps where predicted.
  We have done this on an A100 (measured knee $3.18$ against $\sim 3$ predicted, with
  throughput falling from $\times 1.07$ at four riders to $\times 0.94$ at five);
  the question is whether it lands at $419$ or $375$\,GB/s of effective bandwidth.
  Two hours.
\item \textbf{The goodput experiment of \S\ref{sec:eval:goodput}}, which is the one
  our own substrate cannot run because its link admits too few riders for
  break-even. One day.
\end{enumerate}

We would also publish whatever came back, including a refutation. The coefficients,
the ablation harness, the purity gate and the analysis scripts are all in the
artifact, so a negative result is as reportable as a positive one --- and on this
project's record, four of nine optimisations we attempted had a mechanism that
fired and a benefit that did not.
